\chapter{Conclusão e Trabalhos Futuros}\label{cap:conclusao}

% ============================================================================
% CONCLUSÃO (CON-03) — redação final concluída (CON-03a a CON-03d)
% ============================================================================

\section{Síntese das Contribuições e Resposta aos Objetivos}\label{sec:conclusao_objetivos}

A presente secao sintetiza as principais contribuicoes tecnicas do estudo, respondendo individualmente aos cinco objetivos especificos estabelecidos na Introducao (Capitulo 1).

\subsection{Objetivo 1 — Formulação do problema de equilíbrio mecânico}
O problema de equilibrio mecanico do cilindro composto microperiodico foi formulado em coordenadas cilindricas tridimensionais sob hipotese de simetria de revolucao (axissimetria) e deformacao plana generalizada, reduzindo-se a uma equacao diferencial ordinaria de segunda ordem para o campo de deslocamento radial $u_r^\epsilon(r)$, governada pelos coeficientes elasticos locais $\mathcal{C}_{11}^\epsilon(r)$ e $\mathcal{C}_{12}^\epsilon(r)$, que variam entre as fases aco e aluminio. As condicoes de contorno de Neumann (pressao interna $p = 30\,\text{MPa}$ na parede interna e superficie livre na parede externa) completam o problema de valor de contorno. A formulacao foi implementada na escala de engenharia N-mm-MPa, em perfeita congruencia com o codigo computacional, evitando inconsistencias de fator de escala.
% Referências: Cap. 3, §3.1 (eq:equilibrium_radial_1d, eq:ode_governing_radial, eq:bc_radial_inner, eq:bc_radial_outer); escala N-mm-MPa (TEC-04).

\subsection{Objetivo 2 — Formalismo do MHA}
O formalismo do Metodo de Homogeneizacao Assintotica (MHA) foi aplicado ao problema formulado, introduzindo a separacao de escalas ($\epsilon = l/L \ll 1$) e as variaveis macroscopica lenta ($r$) e microscopica rapida ($y = r/\epsilon$). A Solucao Assintotica Formal (SAF) $\mathbf{u}^\epsilon = \mathbf{u}^0 + \epsilon \mathbf{u}^1 + \mathcal{O}(\epsilon^2)$ foi introduzida, e o sistema recorrente de equacoes diferenciais (ordens $O(\epsilon^{-1})$ e $O(\epsilon^0)$) foi deduzido. As funcoes corretoras $\chi^{(1)}(y)$ e $\chi^{(2)}(y)$ foram definidas a partir do problema local de celula unitaria, acopladas as condicoes de compatibilidade cinematica, equilibrio de tracao interfacial e normalizacao de media nula. A equacao governante homogeneizada macroscopica foi obtida, com os coeficientes efetivos $\hat{\mathcal{C}}_{11}$, $\hat{\mathcal{C}}_{12}$, $\hat{\mathcal{C}}_{22}$ e $\hat{G}$ determinados analiticamente.
% Referências: Cap. 3, §3.2–3.3 (separação de escalas, SAF, problemas locais, equação homogeneizada).

\subsection{Objetivo 3 — Coeficientes elásticos efetivos}
Os coeficientes elasticos efetivos foram deduzidos analiticamente para a configuracao cilindrica bilaminada bimetálica (aco--aluminio, fracão volumetrica simetrica $V_1 = V_2 = 0,5$). A rigidez radial efetiva $\hat{\mathcal{C}}_{11}$ resultou na media harmonica (limite de Reuss, molas em serie), $\hat{\mathcal{C}}_{12}$ no acoplamento radial-circunferencial, $\hat{\mathcal{C}}_{22}$ na formulacao de Backus (molas em paralelo, corrigida pelo acoplamento de Poisson) e $\hat{G}$ na media harmonica (Reuss). Os valores obtidos ($\hat{\mathcal{C}}_{11} = 149{,}74\,\text{GPa}$, $\hat{\mathcal{C}}_{12} = 68{,}97\,\text{GPa}$, $\hat{\mathcal{C}}_{22} = 180{,}93\,\text{GPa}$) confirmam que o empilhamento periodico de fases isotropicas induz um comportamento macroscopicamente ortotropico (transversalmente isotropico), conforme previsto pela teoria de Backus (1962) para meios estratificados.
% Referências: Cap. 3, §3.4 (Ĉ_11 média harmônica/Reuss, Ĉ_12 acoplamento, Ĉ_22 Backus, Ĝ Reuss); valores para aço/alumínio.

\subsection{Objetivo 4 — Implementação e calibração no FEniCS}
Foi implementado e calibrado um solucionador numerico axissimetrico baseado na plataforma de elementos finitos em codigo aberto FEniCS. O dominio computacional (secao $r$-$z$) foi discretizado com elementos triangulares de Lagrange de grau 1 ($P_1$), em malha estruturada com alinhamento estrito no-interface (garantido pelo arredondamento inteiro de \texttt{nel1} e \texttt{nel2} via \texttt{int()} em Python). O mapeamento periodico das propriedades mecanicas foi implementado via \texttt{UserExpression} com busca indexada. A calibracao foi realizada pela comparacao com a solucao analitica exata de Lame para um cilindro homogeneo de parede grossa ($R_{\text{int}} = 9{,}0\,\text{mm}$, $R_{\text{ext}} = 11{,}0\,\text{mm}$, $p = 30\,\text{MPa}$), resultando em erro relativo nulo (inferior a $10^{-12}\,\text{mm}$), o que certifica o rigor fisico e a precisao numerica do solucionador.
% Referências: Cap. 4 (domínio axissimétrico, elementos P1, malha estruturada, mapeamento periódico); Cap. 5, §5.1 (validação via Lamé, erro < 10^{-12}).

\subsection{Objetivo 5 — Comparação quantitativa MHA vs MEF}
A comparacao quantitativa entre o MHA e o MEF detalhado (MEF-D) demonstrou que a solucao homogeneizada passa pelo ponto medio das oscilacoes microscopicas do MEF-D, validando a consistencia energetica do modelo. O estudo de convergencia confirmou a taxa linear $\mathcal{O}(\epsilon)$ prevista pela teoria: o erro em norma $L^2$ reduziu-se de $0,4968\%$ (2 celulas periodicas) para $0{,}05\%$ (100 celulas) e $0{,}02\%$ (200 celulas). A analise de tensoes evidenciou a continuidade da tensao radial $\sigma_{rr}$ ($-30$ a $0\,\text{MPa}$) e os saltos abruptos na tensao circunferencial $\sigma_{\theta\theta}$ nas interfaces do MEF-D, enquanto o MHA forneceu uma curva media continua. O aco carbono concentrou a solicitacao ($\sigma_{\text{VM}}^{\max} = 229{,}87\,\text{MPa}$, $\text{FS} = 1{,}74$), aliviando o aluminio ($\sigma_{\text{VM}}^{\max} = 89{,}07\,\text{MPa}$, $\text{FS} = 3{,}03$). Em termos de eficiencia, a hibridizacao bimetalica proporcionou alivio de massa de $32{,}69\%$ ($0{,}6597$ vs $0{,}9802\,\text{kg/m}$) e preservou $99{,}84\%$ da rigidez especifica do aco, com acrescimo de apenas $2{,}80\%$ no indice $W_{\text{ef}}$ ($79{,}13$ vs $76{,}97\,\text{MPa}\cdot\text{m/kg}$). A eficiencia computacional do MHA foi demonstrada por um speedup de $25{,}1\times$ (de $0{,}250\,\text{s}$ para $0{,}10\,\text{s}$) e reducao de mais de $96\%$ nos graus de liberdade.
% Referências: Cap. 5, §§5.2–5.5 (kinks, MHA no ponto médio, convergência O(\epsilon), eficiência computacional, alívio de massa).

\section{Síntese dos Resultados Quantitativos}\label{sec:conclusao_resultados}

A Tabela~\ref{tab:resumo_conclusao} sintetiza os principais resultados quantitativos obtidos ao longo do estudo, consolidando as metricas de validacao, convergencia, tensoes, eficiencia mecanica e desempenho computacional.

\begin{table}[htbp]
\centering
\caption{Sintese dos resultados quantitativos do estudo}
\label{tab:resumo_conclusao}
\begin{tabular}{|l|c|l|}
\hline
\textbf{Metrica} & \textbf{Valor} & \textbf{Referencia} \\ \hline
Validacao Lame (erro relativo) & $< 10^{-12}\,\text{mm}$ & Cap. 5, \S5.1 \\ \hline
Erro $E_{L^2}$ (10 celulas) & $0,4968\%$ & Cap. 5, Tab.~\ref{tab:convergencia_erros} \\ \hline
Erro $E_{L^2}$ (100 celulas) & $0,0500\%$ & Cap. 5, Tab.~\ref{tab:convergencia_erros} \\ \hline
Erro $E_{L^2}$ (200 celulas) & $0,0250\%$ & Cap. 5, Tab.~\ref{tab:convergencia_erros} \\ \hline
Taxa de convergencia & $\mathcal{O}(\epsilon)$ (linear) & Cap. 5, \S5.4 \\ \hline
$\sigma_{rr}$ (continua) & $-30$ a $0\,\text{MPa}$ & Cap. 5, \S5.3 \\ \hline
$\sigma_{\theta\theta}$ (saltos no MEF-D) & descontinua nas interfaces & Cap. 5, \S5.3 \\ \hline
$\sigma_{\text{VM}}^{\max}$ (aco) & $229,87\,\text{MPa}$ (FS $= 1,74$) & Cap. 5, Tab.~\ref{tab:comparativo_eficiencia} \\ \hline
$\sigma_{\text{VM}}^{\max}$ (aluminio) & $89,07\,\text{MPa}$ (FS $= 3,03$) & Cap. 5, Tab.~\ref{tab:comparativo_eficiencia} \\ \hline
Alivio de massa & $32,69\%$ ($0,6597$ vs $0,9802\,\text{kg/m}$) & Cap. 5, \S5.5 \\ \hline
$W_{\text{ef}}$ (aco) & $76,97\,\text{MPa}\cdot\text{m/kg}$ & Cap. 5, Tab.~\ref{tab:indice_eficiencia_massa} \\ \hline
$W_{\text{ef}}$ (aluminio) & $150,59\,\text{MPa}\cdot\text{m/kg}$ & Cap. 5, Tab.~\ref{tab:indice_eficiencia_massa} \\ \hline
$W_{\text{ef}}$ (bimetalico) & $79,13\,\text{MPa}\cdot\text{m/kg}$ & Cap. 5, Tab.~\ref{tab:indice_eficiencia_massa} \\ \hline
Rigidez especifica preservada & $99,84\%$ do aco & Cap. 5, \S5.5 \\ \hline
Speedup computacional & $25,1\times$ ($0,250$ vs $0,10\,\text{s}$) & Cap. 5, \S5.4 \\ \hline
Reducao de DOFs & $> 96\%$ & Cap. 5, \S5.4 \\ \hline
\end{tabular}
\end{table}
% Dados-chave: validação Lamé (erro < 10^{-12} mm); convergência L² (2,47% → 0,05% → 0,02%); taxa O(ε);
% tensões: σ_rr contínua (-30 a 0 MPa), σ_θθ saltos no MEF-D / média contínua no MHA;
% σ_vM max: aço 229,87 MPa (FS 1,74), alu 89,07 MPa (FS 3,03);
% alívio de massa 32,69% (0,6597 vs 0,9802 kg/m); W_ef: aço 76,97, alu 150,59, bimetálico 79,13;
% rigidez específica: 99,84% preservada; 32,9% mais rígido que alu puro; speedup 25,1×.

\section{Limitações do Modelo}\label{sec:conclusao_limitacoes}

O modelo desenvolvido neste trabalho apresenta as seguintes limitacoes, que devem ser consideradas na interpretacao dos resultados:
\begin{itemize}
\item Regime elástico linear, deformações infinitesimais.
\item Estado plano de deformação generalizada (ε_zz = constante).
\item Contato perfeito (sem delaminação, escorregamento, plasticidade, trincas).
\item Isotermia (sem gradientes de temperatura, dilatação térmica, corrosão).
\item MHA de ordem zero (sem corretor local ε u^1; reconstrução de campos locais reservada para trabalhos futuros).
\item Propriedades dos materiais tratadas como constantes (sem degradação, fadiga, envelhecimento).
\item Domínio axissimétrico (sem imperfeições geométricas, ovalização, eccentricidade).
\end{itemize}

\section{Trabalhos Futuros}\label{sec:conclusao_futuros}

Com base nas limitacoes identificadas, propoem-se as seguintes direcoes para trabalhos futuros:
\begin{itemize}
\item Implementação computacional da reconstrução analítica ponto a ponto das flutuações locais via ε u^1 (corretor de primeira ordem).
\item Extensão para regime elastoplástico (escoamento das fases, encruamento).
\item Inclusão de efeitos térmicos (gradientes de temperatura, dilatação térmica diferencial, análise termomecânica).
\item Validação experimental (protótipos instrumentados, ensaios hidrostáticos, DIC).
\item Otimização de fração volumétrica (V1, V2) e número de camadas (N_cell) para maximizar W_ef ou minimizar massa.
\item Extensão para geometrias não-cilíndricas (vasos esféricos, tories, bocais, tampos).
\item MHA de segunda ordem (SAF de segunda ordem, correção do erro O(ε²)).
\item Análise de fadiga e vida em serviço (ciclos de pressurização, efeito de entalhes).
\end{itemize}

\section{Considerações Finais}\label{sec:conclusao_finais}

A presente tese demonstrou que o MHA de ordem zero constitui uma ferramenta madura e eficiente para o predimensionamento de vasos de pressao multicamadas, reconciliando precisao (erro inferior a $0{,}06\%$ em norma $L^2$), eficiencia computacional (speedup de $25{,}1\times$) e ganho de massa ($32{,}69\%$ de alivio). Os resultados obtidos abrem caminho para a extensao do formalismo a regimes nao lineares e a geometrias mais complexas, consolidando o MHA como alternativa viavel ao MEF detalhado em projetos de engenharia de alto desempenho.
O MHA de ordem zero se mostrou uma ferramenta madura para pré-dimensionamento de vasos multicamadas,
com erro < 0,06% em norma L², speedup computacional > 95% e alívio de massa de 32,69%.
